\documentclass{simpledoc}


\newcommand{\soft}{\textsc{d-Tracker}}
\newcommand{\gui}{\screen{tracker-gui}}
\newcommand{\cli}{\screen{tracker}}
\newcommand{\var}[1]{{\color{myblue}$\langle$#1$\rangle$}}
\newcommand{\opt}[1]{\color{myblue}$[$#1$]$}
\newcommand{\file}[1]{\textit{\color{myblue}#1}}
\newcommand{\screen}[1]{\texttt{\color{myblue}#1}}
\newcommand{\scr}[1]{\texttt{\color{myblue}#1}}
\newcommand{\key}[1]{\fbox{\color{myblue}#1}}

% Placeholder for a screenshot that will be inserted later.
% The argument is a short description of what the figure should show.
\newcommand{\screenshot}[2]{%
  \begin{figure}[H]
  \centering
  \fbox{\parbox[c][4.5cm][c]{0.8\textwidth}{\centering\color{myblue}%
    \textit{[ screenshot placeholder ]}\\[4pt]#2}}
  \caption{#1}
  \end{figure}}

\title{\textsc{Tracker-gui} -- Documentation}
\author{
\textbf{Vincent Richefeu} $\langle$Vincent.Richefeu@3sr-grenoble.fr$\rangle$\\
\textbf{Ga\"{e}l Combe} $\langle$Gael.Combe@3sr-grenoble.fr$\rangle$\\
Laboratoire 3SR, Universit\'{e} Grenoble Alpes, France
}
\date{Version 2026}

\pagestyle{myfootings}
\markboth{Tracker-gui, user guide}%
         {Tracker-gui, user guide}

%%%%%%%%%%%%%%%%%%%%%%%%%%%%%%%%%%%
%%%%%%%%%%%%%%%%%%%%%%%%%%%%%%%%%%%
\begin{document}

\maketitle

\begin{abstract}
This document describes \gui, the graphical companion application of \soft.
It explains how to launch the interface, how to navigate a set of digital image correlation (DIC)
results, how to colour and inspect the tracked points, and how to use it to run procedures and to
prepare manual corrections.
This documentation is a work in progress and will be extended as the application evolves; it
nevertheless aims to be complete for the current version.
\end{abstract}

\tableofcontents

\newpage

% ============================================================
\section{Introduction}

\soft\ is dedicated to 2D digital image correlation. The reference tool of the suite is the
command-line application \cli, fully described in the companion document \emph{Tracker --
Documentation}. \textbf{The CLI is the program that does the actual work}: it reads a command file
(\file{.trk}), runs the selected procedure (particle tracking, distortion correction, sub-pixel
centring, pattern quality, gray-level analysis, visualization, post-processing), and writes the
result files.

\gui\ is an \emph{assistant} built on top of the same engine. It does not replace the CLI and it does
not implement any new correlation algorithm of its own: both binaries are linked against the same
static library \screen{libtracker.a}, so a procedure launched from the GUI behaves exactly as if it
had been launched from \cli. The role of the GUI is to:

\begin{itemize}
  \item \textbf{visualize} the tracked points and the DIC results on top of the images,
  \item \textbf{inspect and tune} parameters (search zones, patterns, colouring of fields),
  \item \textbf{run} a procedure on the current command file and reload the results,
  \item \textbf{assist manual corrections} of grains that were poorly tracked.
\end{itemize}

\attention
{\color{red} This code is designed for academic research. To be used with your own responsibility.}

\screenshot{General view of \gui: the 2D viewport on the left and the control panel on the right.}
{Full window with a set of grains coloured by NCC over a speckle image, the right-hand control panel
visible, and the floating information card at the top left.}

% ============================================================
\section{Building and launching}

\gui\ is built together with \cli. From the \file{src} folder, type:

\begin{verbatim}
$ make
\end{verbatim}

\noindent The \file{Makefile} downloads the required single-header libraries
(\screen{raygui.h}, \screen{rini.h}) and the embedded \textsc{Open Sans} font, then produces two
binaries: \cli\ (the CLI) and \gui\ (the graphical assistant). The GUI is based on the
\screen{raylib}/\screen{raygui} libraries.

The application is launched with a command file as its single argument, exactly the same
\file{.trk} file that \cli\ uses:

\begin{verbatim}
$ tracker-gui <command_file>
\end{verbatim}

\noindent The command file is parsed at start-up (through the same \screen{init()} routine as the CLI),
so all parameters (image names, DIC names, reference image, search zones, pattern shape, procedure,
etc.) are read from it. It is therefore convenient to run \gui\ from the working directory that
contains the images and the \file{dic\_out\_*.txt} files.

\attention Running \gui\ with no argument (or more than one) prints a usage message and exits.

% ============================================================
\section{The interface at a glance}

The window is split into two regions:

\begin{itemize}
  \item a large \textbf{2D viewport} on the left, where the images, grains, patterns, search zones and
        positions are drawn;
  \item a \textbf{control panel} on the right, organized in labelled sections\\ (\scr{NAVIGATION},
        \scr{COLORING}, \scr{POSITIONS}, \scr{BACKGROUND IMAGE},\\ \scr{COMPUTATION}, \scr{CORRECTIONS},
        \scr{SETTINGS}).
\end{itemize}

Floating rounded ``cards'' overlay the viewport with contextual information:

\begin{itemize}
  \item a card at the top left shows the current DIC (or image) number and the number of grains;
  \item when a grain is selected, an \textbf{inspector} card lists all its quantities (see
        Section~\ref{sec:inspector});
  \item when search zones are shown and a grain is selected, a legend card at the bottom left lists the
        normal, rescue and super-rescue zone dimensions;
  \item a vertical \textbf{colour bar} appears at the top right when a field (other than \scr{grey}) is
        used for colouring.
\end{itemize}

Most actions are available both as a button in the panel and as a keyboard shortcut; the shortcut
letter is recalled in parentheses on the corresponding button.

% ============================================================
\section{Navigation}
\label{sec:navigation}

The viewport uses a 2D camera with pan and zoom. The coordinate system is the full-resolution pixel
system of the reference image, with the $y$ axis pointing downward (top of the image is $y=0$), as in
the CLI.

\begin{center}
\begin{tabular}{lp{7cm}}
\hline
\textbf{Action} & \textbf{Effect} \\
\hline
Mouse wheel              & Zoom in/out, centred on the cursor. \\
Right or middle drag     & Pan the view. \\
Left click               & Select the grain under the cursor (outside correction mode). \\
\key{F}                  & Fit the view to all grains. \\
\key{$\leftarrow$} / \key{$\rightarrow$} & Go to the previous / next DIC file (step \scr{iinc}). \\
\hline
\end{tabular}
\end{center}

The \scr{NAVIGATION} section of the panel offers the equivalent \screen{Prev}, \screen{Next} and
\screen{Fit view} buttons. Stepping moves by \scr{iinc} (the increment read from the command file), or
by 1 if \scr{iinc} is not set.

% ============================================================
\section{Keyboard shortcuts}
\label{sec:shortcuts}

Letter shortcuts are read as \emph{typed characters} rather than physical key codes, so they work
identically on QWERTY and AZERTY keyboards.

\begin{center}
\begin{tabular}{lp{8.5cm}}
\hline
\textbf{Key} & \textbf{Action} \\
\hline
\key{F} & Fit the view to all grains. \\
\key{G} & Show / hide the grains (when hidden, grains are drawn as fixed-size screen crosses). \\
\key{P} & Show / hide the correlation pattern of the grains. \\
\key{Z} & Show / hide the search zones. \\
\key{R} & Show / hide the reference / previous / current positions and the rotation indicator. \\
\key{B} & Show / hide the background image. \\
\key{T} & Open the read-only command-file (\file{.trk}) viewer. \\
\key{C} & Enter / leave the manual correction mode for the selected grain. \\
\key{Enter} & In correction mode, register the candidate pose as a correction. \\
\key{Esc} & Close the \file{.trk} viewer if open; otherwise leave correction mode; otherwise quit. \\
\key{$\leftarrow$}/\key{$\rightarrow$} & Previous / next DIC file. \\
\hline
\end{tabular}
\end{center}

% ============================================================
\section{Colouring the tracked points}
\label{sec:coloring}

The \scr{COLORING} section selects which field is mapped to the colour of each grain. The choice is
made through a grid of toggle buttons (two columns). The available fields are:

\begin{center}
\begin{tabular}{lp{8cm}}
\hline
\textbf{Field} & \textbf{Meaning} \\
\hline
\scr{grey}        & No field: grains are drawn in a neutral colour (no colour bar). \\
\scr{NCC}         & Best NCC of the first search. \\
\scr{NCC\_rescue} & Best NCC after rescue. \\
\scr{NCC\_subpix} & NCC after sub-pixel refinement. \\
\scr{dx}, \scr{dy}     & Cumulated displacement components from the reference (pixels). \\
\scr{drot}        & Cumulated rotation from the reference (radians). \\
\scr{upix}, \scr{vpix} & Displacement increment from the previous image (pixels). \\
\scr{rot\_inc}    & Rotation increment from the previous image (radians). \\
\hline
\end{tabular}
\end{center}

The three NCC fields are mapped over the range $[\scr{NCC\_min}, 1]$, where the lower bound is set by
a slider (\scr{ncc\_min}). The other fields use an automatic min/max range computed over the visible
grains; the bounds are displayed below the grid. A vertical colour bar is drawn at the top right of
the viewport for every field except \scr{grey}.

\marginlabel{\key{G} \scr{Show / Hide grains}}
When grains are shown, they are drawn as filled disks (the fill opacity is adjustable with a slider).
When hidden, each grain is drawn instead as a small fixed-size cross, still coloured by the selected
field, which keeps the display readable at any zoom level and over dense point clouds.

\screenshot{Grains coloured by a field, with the vertical colour bar.}
{The viewport with grains coloured by \scr{NCC}, the colour bar visible at the top right, and the
fill-opacity slider set to a partially transparent value.}

% ============================================================
\section{Reference, previous and current positions}
\label{sec:positions}

\marginlabel{\key{R} \scr{Positions}}
The \scr{POSITIONS} section (shortcut \key{R}) overlays, for each grain (or only for the selected one
if \screen{Selected grain only} is enabled), three positions linked by a trace:

\begin{itemize}
  \item the \textbf{reference} position (violet), $(\scr{refcoord\_xpix}, \scr{refcoord\_ypix})$;
  \item the \textbf{previous} position (orange), i.e.\ the current position minus the last increment
        $(\scr{upix}, \scr{vpix})$;
  \item the \textbf{current} position (green), i.e.\ the reference plus the cumulated displacement
        $(\scr{dx}, \scr{dy})$.
\end{itemize}

The reference and previous positions are marked with small circles. On the \textbf{current} position,
a rotation indicator is drawn: a larger green circle with two crossing markers, one at the reference
orientation \scr{refrot} (violet, un-rotated) and one at the current orientation
$\scr{refrot}+\scr{drot}$ (green, rotated). A filled circular sector between the two visualizes the
cumulated rotation \scr{drot}. All these markers have a fixed screen size, so they remain visible at
any zoom. The legend colours are recalled in the panel.

\screenshot{Positions and rotation indicator of a grain.}
{A zoomed-in view of one grain showing the violet reference marker, the orange previous marker, the
green current marker with its rotation sector, and the connecting trace.}

% ============================================================
\section{Pattern and search zones}

\marginlabel{\key{P} \scr{Show / Hide pattern}}
The correlation pattern is the set of pixels of the reference image used for matching. It is not
stored in the DIC files: it is rebuilt by the engine from the command file (the \scr{pattern} shape)
into the internal grain structure, and \gui\ reads it directly. When the true pattern is not
available for a grain, a fixed-size square fallback is drawn instead.

\marginlabel{\key{Z} \scr{Show / Hide search zones}}
The search zones (normal, rescue, super-rescue) are the pixel-resolution displacement windows defined
in the command file. When shown and a grain is selected, the dimensions and rotation parameters of the
three zones are listed in a legend card (colour-coded green / orange / red) at the bottom left.

% ============================================================
\section{Selected-grain inspector}
\label{sec:inspector}

Left-clicking a grain selects it. The selection is persistent across navigation and a dedicated
inspector card lists every quantity of the grain:

\begin{itemize}
  \item \scr{index} and, for results read from a DIC file, the corresponding \scr{file line} and file
        name (\file{dic\_out\_\textit{n}.txt});
  \item \scr{refcoord} (reference pixel position), \scr{radius}, \scr{refrot};
  \item \scr{position} (current position), \scr{dx}, \scr{dy}, \scr{drot};
  \item \scr{upix}, \scr{vpix}, \scr{rot\_inc};
  \item \scr{NCC}, \scr{NCC\_rescue}, \scr{NCC\_subpix};
  \item the number of \scr{pattern} points, when the pattern is available.
\end{itemize}

These are precisely the thirteen columns of the DIC output file (see the CLI documentation), shown for
the selected grain. When the grain radius is not stored, an effective radius is computed for display.

% ============================================================
\section{Background image}

\marginlabel{\key{B} \scr{Show / Hide image}}
The \scr{BACKGROUND IMAGE} section overlays the actual image behind the grains, which is useful to
check the placement of the patterns and the quality of the tracking. The image is read through the
same TIFF / raw readers as the engine. A sub-sampling factor (\scr{image\_div}) reduces memory and
loading time for large images; the thumbnail is scaled back so that the grains, expressed in
full-resolution pixels, remain correctly aligned.

\attention The background follows the current DIC number, so stepping through the images updates the
displayed picture accordingly.

% ============================================================
\section{Command-file viewer}

\marginlabel{\key{T} \scr{View .trk}}
Pressing \key{T} (or the \screen{View .trk} button) opens a read-only overlay showing the full content
of the command file passed on the command line. The text is scrollable with the mouse wheel; \key{T}
or \key{Esc} closes it. This is a convenient reminder of the parameters currently in effect, without
leaving the application.

% ============================================================
\section{Running a procedure}
\label{sec:run}

\marginlabel{\scr{COMPUTATION}}
The \scr{COMPUTATION} section contains a single \screen{Run} button whose label is adapted to the
\scr{procedure} declared in the command file (for example \screen{Run tracking},
\screen{Correct distortion}, \screen{Find centers}, \screen{Pattern quality},\\ \screen{Generate visu},
\screen{Gray analysis}, \screen{Post-process}). Pressing it dispatches to exactly the same procedure
function as \cli\ would call.

Around the run, \gui\ takes a snapshot of the working directory, so that the files \emph{produced or
modified} by the procedure can be detected. A summary banner reports them for a few seconds. The
spatial results are then reloaded automatically:

\begin{enumerate}
  \item the \file{dic\_out\_\textit{n}.txt} of the current number, if present;
  \item otherwise \file{recentered.data} (produced by \scr{find\_subpixel\_centers});
  \item otherwise the grains held in the engine's global state.
\end{enumerate}

If the background was displayed, it is reloaded as well.

\attention The run is currently \emph{synchronous}: the window is unresponsive while a long procedure
is computing. Asynchronous execution is planned (see Section~\ref{sec:limitations}).

% ============================================================
\section{Manual corrections}
\label{sec:corrections}

DIC occasionally mismatches a grain (for example in low-texture regions or near large displacements).
\gui\ helps to fix such grains one by one.

\marginlabel{\key{C} \scr{Correct selected}}
Select the offending grain, then press \key{C} (or \screen{Correct selected}) to enter correction
mode. In this mode you manually place the correct pose of the grain on the deformed image:

\begin{center}
\begin{tabular}{lp{8cm}}
\hline
\textbf{Action} & \textbf{Effect} \\
\hline
Left drag        & Translate the candidate pose ($dx$, $dy$). \\
\key{Shift}+wheel & Rotate the candidate pose ($d\theta$, about $0.5^{\circ}$ per notch). \\
\key{Enter}      & Register the candidate as a correction. \\
\key{C} / \key{Esc} & Leave correction mode. \\
\hline
\end{tabular}
\end{center}

A magenta preview shows the candidate pose while you adjust it. Registered corrections are collected
in a list (de-duplicated per (DIC, grain) pair) and shown in the \scr{CORRECTIONS} section. The
\screen{Save file} button writes them to \file{dic\_corrections.txt}, one correction per line, in the
format:

\begin{verbatim}
dic  grain  dx  dy  drot
\end{verbatim}

\noindent where \scr{drot} is in radians. The \screen{Clear} button empties the list.

\attention \textbf{Important.} In the current version, saving a correction only \emph{records a
request} in \file{dic\_corrections.txt}. It does \emph{not} re-run the correlation and does
\emph{not} modify the \file{dic\_out\_*.txt} files. Re-applying the recorded corrections to the DIC
results (re-running a sub-pixel follow from the corrected pose) is the planned next step
(Section~\ref{sec:limitations}).

\screenshot{Manual correction of a grain.}
{Correction mode active on a selected grain, with the magenta candidate preview offset from the
green current position, and the \scr{CORRECTIONS} list showing one registered entry.}

% ============================================================
\section{Persistent settings}
\label{sec:settings}

\gui\ stores cosmetic and display preferences in an INI file named \file{tracker-gui.ini} in the
current working directory (handled through the \screen{rini} single-header library). The file is read
at start-up, before the window is created (so that the window size is honoured), and written by the
\screen{Save settings} button in the \scr{SETTINGS} section.

\begin{center}
\begin{tabular}{lp{6cm}}
\hline
\textbf{Key} & \textbf{Meaning} \\
\hline
\scr{window\_\opt{width|height}}          & Window size in pixels. \\
\scr{widget\_font\_size}                  & Font size of the widgets. \\
\scr{widget\_button\_height}              & Height of buttons and toggles. \\
\scr{widget\_slider\_height}              & Height of sliders. \\
\scr{widget\_label\_height}               & Vertical step between text lines. \\
\scr{marker\_radius}                      & Base size of the screen markers. \\
\scr{marker\_rotation\_radius}            & Size of the rotation indicator. \\
\scr{marker\_fallback\_radius}            & Default grain radius when none is stored. \\
\scr{image\_div}                          & Sub-sampling factor of the background image. \\
\scr{color\_field}                        & Index of the colouring field. \\
\scr{show\_*}                             & Display toggles. \\
\scr{ncc\_min}                            & Lower bound of the NCC colour range $[0,1]$. \\
\scr{fill\_opacity}                       & Opacity of the grain fill $[0,1]$. \\
\hline
\end{tabular}
\end{center}

% ============================================================
\section{Limitations and roadmap}
\label{sec:limitations}

\gui\ is under active development. Known limitations and planned features include:

\begin{itemize}
  \item \textbf{Apply manual corrections.} Re-run a sub-pixel follow from a corrected pose and update
        the \file{dic\_out\_*.txt} files (with a backup), so that corrections recorded in
        \file{dic\_corrections.txt} actually take effect.
  \item \textbf{Display produced images.} Show the images generated by\\ \scr{visu\_process}
        (\file{Visu\_*.tiff}) directly in the viewport.
  \item \textbf{Asynchronous runs.} Keep the interface responsive while a procedure is computing.
\end{itemize}

As the engine and the GUI evolve, this documentation will be extended accordingly. For everything
related to the procedures themselves, the parameters of the command file and the output formats, refer
to the companion document \emph{Tracker -- Documentation}.

\end{document}
